% ID: SEL-D0046-OB-E01
% ESTADO: validada. S047.
\textbf{a)} Sobre $y=0$, las otras restricciones dejan $-2\le x\le 5/2$. La frontera inferior izquierda es $y=-4-2x$ y corta a $y=x+5$ en $(-3,2)$. Las dos cotas superiores, $x+5$ y $10-4x$, se igualan en $(1,6)$. Comprobando las cuatro restricciones, los vértices en orden son $(-3,2),(-2,0),(5/2,0),(1,6)$.
\begin{center}
\setlength{\unitlength}{1mm}
\begin{picture}(103,78)(-53,-8)
\qbezier(-38,0)(0.5,0)(39,0)
\qbezier(0,-7)(0,31)(0,69)
\put(40,0){$x$}
\put(2,67){$y$}
{
\qbezier(-30,20)(-25,10)(-20,0)
\qbezier(-20,0)(2.5,0)(25,0)
\qbezier(25,0)(17.5,30)(10,60)
\qbezier(10,60)(-10,40)(-30,20)
}
\put(-52,21){\scriptsize $(-3,2)$}
\put(-29,-6){\scriptsize $(-2,0)$}
\put(16,-6){\scriptsize $(5/2,0)$}
\put(12,60){\scriptsize $(1,6)$}
\end{picture}
\end{center}
\textbf{b)} Evaluar la función en los vértices da, respectivamente, $-2,-2,5/2,4$. El máximo es $\boxed{4}$ sólo en $(1,6)$. El mínimo es $\boxed{-2}$ en todo el segmento entre $(-3,2)$ y $(-2,0)$: la igualdad $f=-2$ equivale a $2x+y=-4$, que define ese lado del recinto.
